% ===== PRÉAMBULE CANONIQUE — à coller VERBATIM en tête de CHAQUE .tex =====
\documentclass[11pt,a4paper]{article}
\usepackage[utf8]{inputenc}\usepackage[T1]{fontenc}\usepackage{lmodern}
\usepackage[french]{babel}
\usepackage{amsmath,amssymb,mathtools}\usepackage{icomma}
\usepackage{siunitx}\sisetup{locale=FR}  % \qty{}{}, \si{}, \ang{}
\usepackage{xcolor}\usepackage{colortbl}
\usepackage{tikz}\usepackage{pgfplots}\pgfplotsset{compat=1.18}
\usepackage[siunitx,europeanresistors]{circuitikz} % circuits (élec.)
\usepackage[most]{tcolorbox}\usepackage{enumitem}\usepackage{array,booktabs}
\usepackage[a4paper,margin=2.2cm]{geometry}
\usetikzlibrary{arrows.meta,calc,angles,quotes,positioning,patterns,%
  backgrounds,shapes.geometric,decorations.pathmorphing,decorations.markings,intersections}
\definecolor{primary}{RGB}{222,49,99}\definecolor{primarydark}{RGB}{150,22,55}
\definecolor{defcol}{RGB}{0,114,178}\definecolor{methcol}{RGB}{52,92,148}
\definecolor{excol}{RGB}{0,135,90}\definecolor{attncol}{RGB}{200,40,55}
\tcbset{boxbase/.style={enhanced,breakable,boxrule=0.9pt,arc=2.5pt,
  left=3mm,right=3mm,top=2.3mm,bottom=2.3mm,fonttitle=\bfseries,coltitle=white}}
% --- boîtes TUTORIEL (noms EXACTS, ne pas renommer) ---
\newtcolorbox{cle}[1][]{boxbase,colback=defcol!6,colframe=defcol,
  title={\textsc{À retenir}\ifx\relax#1\relax\relax\else\textmd{ : \itshape #1}\fi}}
\newtcolorbox{methode}[1][]{boxbase,colback=methcol!7,colframe=methcol,sharp corners,
  title={\textsc{Méthode}\ifx\relax#1\relax\relax\else\textmd{ --- \itshape #1}\fi}}
\newtcolorbox{exemplebox}[1][]{boxbase,colback=excol!5,colframe=excol,
  title={\textsc{Exemple}\ifx\relax#1\relax\relax\else\textmd{ : \itshape #1}\fi}}
\newtcolorbox{piege}[1][]{boxbase,colback=attncol!6,colframe=attncol,
  title={\textsc{Piège}\ifx\relax#1\relax\relax\else\textmd{ : \itshape #1}\fi}}
\newtcolorbox{remarque}[1][]{boxbase,colback=black!4,colframe=black!45,
  title={\textsc{Remarque}\ifx\relax#1\relax\relax\else\textmd{ : \itshape #1}\fi}}
% --- boîtes EXERCICE / CORRIGÉ (noms EXACTS, auto-comptées) ---
\newtcolorbox[auto counter]{exo}[1][]{boxbase,colback=primary!4,colframe=primary,
  title={\textsc{Exercice}~\thetcbcounter\ifx\relax#1\relax\relax\else\textmd{ --- \itshape #1}\fi}}
\newtcolorbox[auto counter]{corrige}[1][]{boxbase,colback=excol!4,colframe=excol,
  title={\textsc{Corrigé}~\thetcbcounter\ifx\relax#1\relax\relax\else\textmd{ --- \itshape #1}\fi}}
\newcommand{\vv}[1]{\vec{#1}}\newcommand{\ds}{\displaystyle}
\newcommand{\R}{\mathbb{R}}\newcommand{\N}{\mathbb{N}}\newcommand{\Z}{\mathbb{Z}}
% (un \usepackage{fancyhdr}+en-tête est possible ; il est ignoré côté web)
% =========================================================================
\begin{document}
\sisetup{per-mode=symbol}

\begin{exo}[projeter un vecteur vitesse]
Décompose la vitesse sur les axes ($v_x=v\cos\theta$, $v_y=v\sin\theta$).
\begin{enumerate}[label=\alph*)]
  \item $v=\qty{10}{\metre\per\second}$, $\theta=\ang{30}$ avec l'horizontale.
  \item $v=\qty{10}{\metre\per\second}$, $\theta=\ang{60}$.
\end{enumerate}
\end{exo}

\begin{exo}[reconstituer la norme et l'angle]
On connaît les composantes ; retrouve la norme $v=\sqrt{v_x^2+v_y^2}$ et l'angle $\theta$.
\begin{enumerate}[label=\alph*)]
  \item $v_x=\qty{6}{\metre\per\second}$, $v_y=\qty{8}{\metre\per\second}$.
  \item $v_x=\qty{5}{\metre\per\second}$, $v_y=\qty{12}{\metre\per\second}$.
\end{enumerate}
\end{exo}

\begin{exo}[composition à angle droit]
Un bateau avance à \qty{8}{\metre\per\second} perpendiculairement à la rive ; le courant, parallèle
à la rive, vaut \qty{6}{\metre\per\second}.
\begin{enumerate}[label=\alph*)]
  \item Quelle est la vitesse du bateau par rapport à la rive (le sol) ?
  \item De quel angle sa trajectoire est-elle déviée par rapport à sa direction propre ?
\end{enumerate}
\end{exo}

\begin{exo}[composition sur une même droite]
Sur un tapis roulant d'aéroport qui avance à \qty{0,8}{\metre\per\second}, un voyageur marche à
\qty{1,5}{\metre\per\second} par rapport au tapis.
\begin{enumerate}[label=\alph*)]
  \item Quelle est sa vitesse par rapport au sol s'il marche dans le sens du tapis ?
  \item Et s'il marche à contre-sens (vers l'arrière) ?
\end{enumerate}
\end{exo}

\begin{exo}[norme constante ou vecteur constant ?]
Réponds par vrai ou faux et justifie brièvement.
\begin{enumerate}[label=\alph*)]
  \item Un mobile dont la norme de la vitesse est constante a forcément un vecteur vitesse constant.
  \item Un mobile en mouvement rectiligne uniforme (MRU) a un vecteur vitesse constant.
\end{enumerate}
\end{exo}

\end{document}
